% DMA Session 10: Subqueries, Views & Transaktionen
% 180-Minuten-Block (Vorlesung + Übung interwoven)

\documentclass[usenames,dvipsnames,10pt,aspectratio=169]{beamer}
\usepackage[T1]{fontenc}
\usepackage[utf8]{inputenc}
\usepackage{verbatim}

\usetheme{ims}

\usepackage{booktabs}
\usepackage{multicol}
\usepackage{listings}
\usepackage[table]{xcolor}
\usepackage{graphicx}
\usepackage{tikz}
\usetikzlibrary{shapes.geometric, arrows.meta, positioning, calc, fit, backgrounds, matrix}
\usepackage{pifont}

% TikZ styles
\tikzset{
    tablebox/.style={rectangle, draw=IMSBlue, fill=IMSBlue!5, thick,
        minimum width=4cm, minimum height=0.8cm, font=\ttfamily\small},
    tableheader/.style={rectangle, draw=IMSBlue, fill=IMSBlue!20, thick,
        minimum width=4cm, minimum height=0.8cm, font=\ttfamily\small\bfseries},
    roadmapbox/.style={rectangle, rounded corners, minimum width=2.5cm, minimum height=0.8cm,
        text centered, draw=IMSBlue, fill=IMSBlue!10, font=\small\bfseries},
    roadmaparrow/.style={-{Stealth[length=3mm]}, thick, IMSBlue},
    sqlbox/.style={rectangle, rounded corners, draw=IMSOrange, fill=IMSOrange!10,
        minimum width=2cm, minimum height=0.6cm, font=\small},
    sqlarrow/.style={-{Stealth[length=2.5mm]}, thick, IMSOrange},
    acidbox/.style={rectangle, rounded corners, minimum width=2.2cm, minimum height=1.2cm,
        text centered, align=center, draw=IMSBlue, thick, font=\small\bfseries, fill=IMSBlue!10},
    problembox/.style={rectangle, rounded corners, minimum width=3cm, minimum height=1cm,
        text centered, draw=red!70, thick, font=\small, fill=red!10},
    solutionbox/.style={rectangle, rounded corners, minimum width=3cm, minimum height=1cm,
        text centered, draw=green!70!black, thick, font=\small, fill=green!10},
    querybox/.style={rectangle, rounded corners, draw=IMSBlue, fill=gray!10,
        minimum width=5cm, minimum height=1.5cm, font=\ttfamily\footnotesize},
}

% SQL Listing Style
\lstdefinestyle{sql}{
    language=SQL,
    basicstyle=\ttfamily\footnotesize,
    keywordstyle=\color{IMSBlue}\bfseries,
    stringstyle=\color{IMSOrange},
    commentstyle=\color{gray}\itshape,
    showstringspaces=false,
    breaklines=true,
    frame=single,
    backgroundcolor=\color{gray!10},
    morekeywords={SERIAL, BOOLEAN, TEXT, REFERENCES, CONSTRAINT, CASCADE, RESTRICT, PRIMARY, FOREIGN, KEY, UNIQUE, CHECK, DEFAULT, AUTO_INCREMENT, INT, VARCHAR, DATE, DECIMAL, NOT, NULL, WITH, AS, BEGIN, COMMIT, ROLLBACK, TRANSACTION, VIEW, CREATE, DROP, REPLACE, EXISTS, ANY, ALL, SOME, IN, RECURSIVE, TEMPORARY, TEMP, MATERIALIZED},
    literate={ü}{{\"u}}1 {ä}{{\"a}}1 {ö}{{\"o}}1 {Ü}{{\"U}}1 {Ä}{{\"A}}1 {Ö}{{\"O}}1 {ß}{{\ss}}1
}

\lstset{style=sql}

\newcommand{\cmark}{\textcolor{green!70!black}{\ding{51}}}
\newcommand{\xmark}{\textcolor{red}{\ding{55}}}
\newcommand{\pk}[1]{\underline{\texttt{#1}}}
\newcommand{\fk}[1]{\texttt{\textcolor{IMSOrange}{#1}}}

% ===== CLICKABLE AGENDA WITH PROGRESS INDICATOR =====
\usepackage{hyperref}
\hypersetup{colorlinks=false, pdfborder={0 0 0}}

% Phase counter for progress tracking
\newcounter{currentphase}
\setcounter{currentphase}{0}

% Clickable agenda item
\newcommand{\agendaitem}[3]{%
    \ifnum#1=#2
        \textcolor{IMSOrange}{$\blacktriangleright$ \textbf{\hyperlink{phase#2}{#3}}}%
    \else
        \textcolor{gray!70}{\phantom{$\blacktriangleright$} \hyperlink{phase#2}{#3}}%
    \fi\\[0.3em]%
}

% Progress dots for footline (clickable)
\newcommand{\progressdots}{%
    \begin{tikzpicture}[baseline=-0.5ex]
        \foreach \i in {1,...,7} {
            \ifnum\value{currentphase}=\i
                \node[circle, fill=IMSOrange, minimum size=0.24cm, inner sep=0pt] at (\i*0.4,0) {\hyperlink{phase\i}{\phantom{oo}}};
            \else
                \ifnum\value{currentphase}>\i
                    \node[circle, fill=IMSBlue!60, minimum size=0.2cm, inner sep=0pt] at (\i*0.4,0) {\hyperlink{phase\i}{\phantom{oo}}};
                \else
                    \node[circle, draw=gray!50, minimum size=0.2cm, inner sep=0pt] at (\i*0.4,0) {\hyperlink{phase\i}{\phantom{oo}}};
                \fi
            \fi
        }
    \end{tikzpicture}%
}

% Add progress indicator to footline
\setbeamertemplate{footline}{%
    \leavevmode%
    \hbox{%
        \begin{beamercolorbox}[wd=.33\paperwidth,ht=2.5ex,dp=1ex,left]{author in head/foot}%
            \usebeamerfont{author in head/foot}\hspace*{2ex}\insertshortauthor
        \end{beamercolorbox}%
        \begin{beamercolorbox}[wd=.34\paperwidth,ht=2.5ex,dp=1ex,center]{title in head/foot}%
            \progressdots
        \end{beamercolorbox}%
        \begin{beamercolorbox}[wd=.33\paperwidth,ht=2.5ex,dp=1ex,right]{date in head/foot}%
            \usebeamerfont{date in head/foot}\insertframenumber{} / \inserttotalframenumber\hspace*{2ex}
        \end{beamercolorbox}%
    }%
    \vskip0pt%
}

% Show agenda (with hyperlink target)
\newcommand{\showagenda}[1]{%
    \setcounter{currentphase}{#1}%
    \hypertarget{phase#1}{}%
    \begin{frame}{Agenda}
    \vfill
    \begin{center}
    \begin{minipage}{0.7\textwidth}
    \large
    \agendaitem{#1}{1}{1 ~ Rückblick \& Motivation}
    \agendaitem{#1}{2}{2 ~ Subqueries: Grundlagen}
    \agendaitem{#1}{3}{3 ~ Common Table Expressions (CTEs)}
    \agendaitem{#1}{4}{4 ~ Views: Virtuelle Tabellen}
    \agendaitem{#1}{5}{5 ~ Transaktionen: ACID}
    \agendaitem{#1}{6}{6 ~ Concurrency-Probleme}
    \agendaitem{#1}{7}{7 ~ Zusammenfassung}
    \end{minipage}
    \end{center}
    \vfill
    \end{frame}
}

%%%%%%%%%%%%%%%%%%%%%%%%%%%%%%%%%%%%%%%%%%%%%%%%%%%%%%%%%%%%%%%%%%%%%%%%%%%%%%%%%%%%%
\title[DMA 10]{Datenmanagement \& -analyse}
\subtitle{Vorlesung 10: Subqueries, Views \& Transaktionen}
\author{Prof. Dr. Christoph M. Flath}
\institute{Lehrstuhl für Wirtschaftsinformatik und Business Analytics\\Julius-Maximilians-Universität Würzburg}
\date{Sommersemester 2026}
%%%%%%%%%%%%%%%%%%%%%%%%%%%%%%%%%%%%%%%%%%%%%%%%%%%%%%%%%%%%%%%%%%%%%%%%%%%%%%%%%%%%%

\begin{document}

% ===== TITLE =====
\begin{frame}[plain]
    \titlepage
\end{frame}

%%%%%%%%%%%%%%%%%%%%%%%%%%%%%%%%%%%%%%%%%%%%%%%%%%%%%%%%%%%%%%%%%%%%%%%%%%%%%%%%%%%%%
\showagenda{1}
%%%%%%%%%%%%%%%%%%%%%%%%%%%%%%%%%%%%%%%%%%%%%%%%%%%%%%%%%%%%%%%%%%%%%%%%%%%%%%%%%%%%%

\begin{frame}{Rückblick: JOINs}
\begin{columns}[T]
\begin{column}{0.5\textwidth}
\textbf{Was wir gelernt haben:}
\begin{itemize}
    \item INNER JOIN: Schnittmenge
    \item LEFT/RIGHT JOIN: Erhalt einer Seite
    \item Self-Joins: Tabelle mit sich selbst
    \item Join-Performance: Indizes nutzen
\end{itemize}
\end{column}

\begin{column}{0.5\textwidth}
\begin{center}
\begin{tikzpicture}[scale=0.7]
    % Venn diagram for INNER JOIN
    \fill[IMSBlue!30] (0,0) circle (1.2);
    \fill[IMSOrange!30] (1.5,0) circle (1.2);
    \begin{scope}
        \clip (0,0) circle (1.2);
        \fill[green!40] (1.5,0) circle (1.2);
    \end{scope}
    \draw[thick, IMSBlue] (0,0) circle (1.2) node[left=0.8cm] {\small A};
    \draw[thick, IMSOrange] (1.5,0) circle (1.2) node[right=0.8cm] {\small B};
    \node at (0.75,-1.8) {\footnotesize INNER JOIN};
\end{tikzpicture}
\end{center}
\end{column}
\end{columns}

\vspace{1em}

\begin{block}{Heute: Komplexere Abfragen strukturieren}
Wie können wir Abfragen \textbf{modular} und \textbf{wiederverwendbar} gestalten?
\end{block}
\end{frame}

\begin{frame}[fragile]{Motivation: Komplexe Abfragen}
\textbf{Problem:} Abfragen werden schnell unübersichtlich

\begin{lstlisting}
SELECT Mannschaft, Punkte,
       (SELECT AVG(Punkte) FROM bundesliga) AS Durchschnitt,
       Punkte - (SELECT AVG(Punkte) FROM bundesliga) AS Differenz
FROM bundesliga
WHERE Punkte > (SELECT AVG(Punkte) FROM bundesliga)
ORDER BY Punkte DESC;
\end{lstlisting}

\vspace{1em}

\begin{alertblock}{Probleme dieser Abfrage}
\begin{itemize}
    \item Durchschnitt wird 3x berechnet (redundant)
    \item Schwer zu lesen und zu warten
    \item Fehleranfällig bei Änderungen
\end{itemize}
\end{alertblock}
\end{frame}

\begin{frame}{Lösungsansätze}
\begin{center}
\begin{tikzpicture}[scale=0.85, transform shape]
    \node[roadmapbox, minimum width=3cm] (sub) at (0,0) {Subqueries};
    \node[roadmapbox, minimum width=3cm] (cte) at (4.5,0) {CTEs (WITH)};
    \node[roadmapbox, minimum width=3cm] (view) at (9,0) {Views};

    \draw[roadmaparrow] (sub) -- (cte) node[midway, above, font=\small] {lesbarer};
    \draw[roadmaparrow] (cte) -- (view) node[midway, above, font=\small] {persistent};

    \node[font=\footnotesize, text width=2.8cm, align=center] at (0,-1.2) {Abfrage in\\Abfrage};
    \node[font=\footnotesize, text width=2.8cm, align=center] at (4.5,-1.2) {Benannte\\Zwischenergebnisse};
    \node[font=\footnotesize, text width=2.8cm, align=center] at (9,-1.2) {Gespeicherte\\Abfrage};
\end{tikzpicture}
\end{center}

\vspace{1em}

\textbf{Und danach:} Transaktionen -- wie Änderungen sicher durchgeführt werden
\end{frame}

%%%%%%%%%%%%%%%%%%%%%%%%%%%%%%%%%%%%%%%%%%%%%%%%%%%%%%%%%%%%%%%%%%%%%%%%%%%%%%%%%%%%%
\showagenda{2}
%%%%%%%%%%%%%%%%%%%%%%%%%%%%%%%%%%%%%%%%%%%%%%%%%%%%%%%%%%%%%%%%%%%%%%%%%%%%%%%%%%%%%

\begin{frame}[fragile]{Subquery-Typen}
\begin{center}
\begin{tikzpicture}
    \node[acidbox, minimum width=3cm, fill=IMSBlue!20] (scalar) at (0,0) {\shortstack{Scalar\\(1 Wert)}};
    \node[acidbox, minimum width=3cm, fill=IMSOrange!20] (column) at (4.5,0) {\shortstack{Column\\(1 Spalte)}};
    \node[acidbox, minimum width=3cm, fill=green!20] (table) at (9,0) {\shortstack{Table\\(Tabelle)}};
\end{tikzpicture}
\end{center}

\vspace{1em}

\begin{columns}[T]
\begin{column}{0.33\textwidth}
\textbf{Scalar Subquery:}
\begin{lstlisting}[basicstyle=\ttfamily\tiny]
SELECT Mannschaft,
  Punkte - (
    SELECT AVG(Punkte)
    FROM bundesliga
  ) AS Diff
FROM bundesliga;
\end{lstlisting}
\end{column}

\begin{column}{0.33\textwidth}
\textbf{Column Subquery:}
\begin{lstlisting}[basicstyle=\ttfamily\tiny]
SELECT Mannschaft
FROM bundesliga
WHERE Mannschaft IN (
  SELECT Mannschaft
  FROM champions
);
\end{lstlisting}
\end{column}

\begin{column}{0.33\textwidth}
\textbf{Table Subquery:}
\begin{lstlisting}[basicstyle=\ttfamily\tiny]
SELECT t.Mannschaft
FROM (
  SELECT Mannschaft,
         Punkte
  FROM bundesliga
  WHERE Punkte > 50
) AS t;
\end{lstlisting}
\end{column}
\end{columns}
\end{frame}

\begin{frame}[fragile]{Scalar Subquery: Ein Wert}
\textbf{Gibt genau einen Wert zurück} -- kann überall stehen, wo ein Wert erwartet wird

\begin{lstlisting}
-- Durchschnittspunkte aller Teams
SELECT Mannschaft, Punkte,
       (SELECT AVG(Punkte) FROM bundesliga) AS Liga_Schnitt
FROM bundesliga
ORDER BY Punkte DESC;
\end{lstlisting}

\vspace{0.5em}

\begin{columns}[T]
\begin{column}{0.5\textwidth}
\textbf{Verwendbar in:}
\begin{itemize}
    \item SELECT (wie oben)
    \item WHERE-Bedingungen
    \item HAVING-Bedingungen
    \item Berechnungen
\end{itemize}
\end{column}

\begin{column}{0.5\textwidth}
\begin{alertblock}{Achtung}
Scalar Subquery muss \textbf{genau 1 Zeile} und \textbf{1 Spalte} liefern!

Sonst: Fehler.
\end{alertblock}
\end{column}
\end{columns}
\end{frame}

\begin{frame}[fragile]{Column Subquery: IN, ANY, ALL}
\textbf{Gibt eine Liste von Werten zurück}

\begin{lstlisting}
-- Teams, die im Pokal-Halbfinale waren
SELECT Mannschaft, Punkte
FROM bundesliga
WHERE Mannschaft IN (
    SELECT team FROM pokal_halbfinale
);
\end{lstlisting}

\vspace{0.5em}

\begin{columns}[T]
\begin{column}{0.5\textwidth}
\textbf{Operatoren:}
\begin{itemize}
    \item \texttt{IN}: Wert in Liste?
    \item \texttt{NOT IN}: Wert nicht in Liste?
    \item \texttt{ANY/SOME}: mind. einer erfüllt
    \item \texttt{ALL}: alle erfüllen
\end{itemize}
\end{column}

\begin{column}{0.5\textwidth}
\begin{lstlisting}[basicstyle=\ttfamily\tiny]
-- Mehr Punkte als IRGENDEIN
-- Team aus Berlin
WHERE Punkte > ANY (
  SELECT Punkte
  FROM bundesliga
  WHERE Mannschaft LIKE '%Berlin%'
);

-- Mehr Punkte als ALLE
-- Teams aus Berlin
WHERE Punkte > ALL (...);
\end{lstlisting}
\end{column}
\end{columns}
\end{frame}

\begin{frame}[fragile]{Korrelierte Subquery}
\textbf{Referenziert die äußere Abfrage} -- wird für jede Zeile neu ausgeführt

\begin{lstlisting}
-- Teams mit überdurchschnittlich vielen Siegen für ihre Punktzahl
SELECT m1.Mannschaft, m1.Punkte, m1.Siege
FROM bundesliga m1
WHERE m1.Siege > (
    SELECT AVG(m2.Siege)
    FROM bundesliga m2
    WHERE m2.Punkte BETWEEN m1.Punkte - 5 AND m1.Punkte + 5
);
\end{lstlisting}

\vspace{0.5em}

\begin{center}
\begin{tikzpicture}
    \node[querybox, minimum width=4cm, align=center] (outer) at (0,0) {Äußere Abfrage\\(m1)};
    \node[querybox, minimum width=4cm, fill=IMSOrange!10, draw=IMSOrange, align=center] (inner) at (6,0) {Innere Abfrage\\(m2)};

    \draw[sqlarrow, <->] (outer) -- (inner) node[midway, above, font=\small] {Referenz auf m1.Punkte};
\end{tikzpicture}
\end{center}

\begin{alertblock}{Performance}
Korrelierte Subqueries können langsam sein -- für jede Zeile wird die innere Abfrage ausgeführt!
\end{alertblock}
\end{frame}

\begin{frame}[fragile]{EXISTS: Existenzprüfung}
\textbf{Prüft, ob Subquery Ergebnisse liefert}

\begin{lstlisting}
-- Teams, die mindestens ein Heimspiel mit 5+ Toren hatten
SELECT DISTINCT m.Mannschaft
FROM bundesliga m
WHERE EXISTS (
    SELECT 1 FROM spiele s
    WHERE s.Heim = m.Mannschaft
    AND s.Heimtore >= 5
);
\end{lstlisting}

\vspace{0.5em}

\begin{columns}[T]
\begin{column}{0.5\textwidth}
\textbf{EXISTS vs. IN:}
\begin{itemize}
    \item EXISTS: Stoppt beim ersten Treffer
    \item IN: Prüft alle Werte
    \item EXISTS oft schneller bei großen Subqueries
\end{itemize}
\end{column}

\begin{column}{0.5\textwidth}
\begin{exampleblock}{Tipp}
\texttt{SELECT 1} in EXISTS -- der Wert ist egal, nur die Existenz zählt.
\end{exampleblock}
\end{column}
\end{columns}
\end{frame}

%%%%%%%%%%%%%%%%%%%%%%%%%%%%%%%%%%%%%%%%%%%%%%%%%%%%%%%%%%%%%%%%%%%%%%%%%%%%%%%%%%%%%
% HANDS-ON 1
{
\setbeamercolor{background canvas}{bg=IMSOrange!15}
\begin{frame}[plain]
\vfill
\begin{center}
{\Huge\color{IMSOrange} Hands-on}\\[1em]
{\Large Subqueries üben}\\[2em]
{\large\ttfamily marimo: 10-subqueries-views-transaktionen.py}\\[1em]
{\normalsize Aufgaben 10.1 -- 10.3}
\end{center}
\vfill
\end{frame}
}

%%%%%%%%%%%%%%%%%%%%%%%%%%%%%%%%%%%%%%%%%%%%%%%%%%%%%%%%%%%%%%%%%%%%%%%%%%%%%%%%%%%%%
\showagenda{3}
%%%%%%%%%%%%%%%%%%%%%%%%%%%%%%%%%%%%%%%%%%%%%%%%%%%%%%%%%%%%%%%%%%%%%%%%%%%%%%%%%%%%%

\begin{frame}[fragile]{Common Table Expressions (CTEs)}
\textbf{Benannte Zwischenergebnisse} -- macht komplexe Abfragen lesbar

\begin{lstlisting}
WITH durchschnitt AS (
    SELECT AVG(Punkte) AS avg_punkte
    FROM bundesliga
)
SELECT Mannschaft, Punkte,
       Punkte - (SELECT avg_punkte FROM durchschnitt) AS Differenz
FROM bundesliga, durchschnitt
WHERE Punkte > durchschnitt.avg_punkte
ORDER BY Punkte DESC;
\end{lstlisting}

\vspace{0.5em}

\begin{block}{Vorteile von CTEs}
\begin{itemize}
    \item Lesbarkeit: Komplexe Logik in benannte Blöcke aufteilen
    \item Wiederverwendung: CTE kann mehrfach referenziert werden
    \item Debugging: Einzelne CTEs separat testen
\end{itemize}
\end{block}
\end{frame}

\begin{frame}[fragile]{CTE Syntax}
\begin{center}
\begin{tikzpicture}
    \node[sqlbox, minimum width=8cm, minimum height=3cm, font=\ttfamily\small, align=left] at (0,0) {
        \textcolor{IMSBlue}{\textbf{WITH}} cte\_name \textcolor{IMSBlue}{\textbf{AS}} (\\
        \quad\quad SELECT ...\\
        ),\\
        weitere\_cte \textcolor{IMSBlue}{\textbf{AS}} (\\
        \quad\quad SELECT ...\\
        )\\
        \textcolor{IMSBlue}{\textbf{SELECT}} ... \textcolor{IMSBlue}{\textbf{FROM}} cte\_name ...
    };
\end{tikzpicture}
\end{center}

\vspace{0.5em}

\begin{itemize}
    \item CTEs werden mit \texttt{WITH} eingeleitet
    \item Mehrere CTEs durch Komma getrennt
    \item Jede CTE hat einen Namen und eine Abfrage
    \item Hauptabfrage folgt nach der letzten CTE
\end{itemize}
\end{frame}

\begin{frame}[fragile]{Mehrere CTEs verketten}
\begin{lstlisting}
WITH
-- Schritt 1: Nur Top-Teams
top_teams AS (
    SELECT Mannschaft, Punkte, Siege
    FROM bundesliga
    WHERE Punkte > 50
),
-- Schritt 2: Statistiken berechnen
team_stats AS (
    SELECT Mannschaft,
           Punkte,
           CAST(Siege AS FLOAT) / NULLIF(Spiele, 0) AS Siegquote
    FROM top_teams
)
-- Hauptabfrage
SELECT * FROM team_stats
ORDER BY Siegquote DESC;
\end{lstlisting}

\vspace{0.5em}

\textbf{Tipp:} Jede CTE kann vorherige CTEs referenzieren!
\end{frame}

\begin{frame}[fragile]{Rekursive CTEs}
\textbf{Für hierarchische Daten} (Org-Charts, Stücklisten, Graphen)

\begin{lstlisting}
WITH RECURSIVE mitarbeiter_hierarchie AS (
    -- Basisfall: CEO (kein Vorgesetzter)
    SELECT id, name, chef_id, 1 AS ebene
    FROM mitarbeiter
    WHERE chef_id IS NULL

    UNION ALL

    -- Rekursion: Mitarbeiter des aktuellen Levels
    SELECT m.id, m.name, m.chef_id, h.ebene + 1
    FROM mitarbeiter m
    JOIN mitarbeiter_hierarchie h ON m.chef_id = h.id
)
SELECT * FROM mitarbeiter_hierarchie;
\end{lstlisting}

\begin{alertblock}{Vorsicht}
Endlosschleifen vermeiden! Immer Abbruchbedingung sicherstellen.
\end{alertblock}
\end{frame}

%%%%%%%%%%%%%%%%%%%%%%%%%%%%%%%%%%%%%%%%%%%%%%%%%%%%%%%%%%%%%%%%%%%%%%%%%%%%%%%%%%%%%
% HANDS-ON 2
{
\setbeamercolor{background canvas}{bg=IMSOrange!15}
\begin{frame}[plain]
\vfill
\begin{center}
{\Huge\color{IMSOrange} Hands-on}\\[1em]
{\Large CTEs strukturieren Abfragen}\\[2em]
{\large\ttfamily marimo: 10-subqueries-views-transaktionen.py}\\[1em]
{\normalsize Aufgaben 10.4 -- 10.5}
\end{center}
\vfill
\end{frame}
}

%%%%%%%%%%%%%%%%%%%%%%%%%%%%%%%%%%%%%%%%%%%%%%%%%%%%%%%%%%%%%%%%%%%%%%%%%%%%%%%%%%%%%
% PAUSE
{
\setbeamercolor{background canvas}{bg=gray!20}
\begin{frame}[plain]
\vfill
\begin{center}
{\Huge ☕ Pause}\\[1em]
{\Large 15 Minuten}
\end{center}
\vfill
\end{frame}
}

%%%%%%%%%%%%%%%%%%%%%%%%%%%%%%%%%%%%%%%%%%%%%%%%%%%%%%%%%%%%%%%%%%%%%%%%%%%%%%%%%%%%%
\showagenda{4}
%%%%%%%%%%%%%%%%%%%%%%%%%%%%%%%%%%%%%%%%%%%%%%%%%%%%%%%%%%%%%%%%%%%%%%%%%%%%%%%%%%%%%

\begin{frame}[fragile]{Views: Virtuelle Tabellen}
\textbf{Gespeicherte Abfrage}, die sich wie eine Tabelle verhält

\begin{lstlisting}
CREATE VIEW top_teams AS
SELECT Mannschaft, Punkte, Tordifferenz
FROM bundesliga
WHERE Punkte > 50;
\end{lstlisting}

\vspace{0.5em}

\begin{lstlisting}
-- Verwendung wie eine normale Tabelle
SELECT * FROM top_teams ORDER BY Punkte DESC;
\end{lstlisting}

\vspace{0.5em}

\begin{columns}[T]
\begin{column}{0.5\textwidth}
\textbf{View ist:}
\begin{itemize}
    \item Keine Datenkopie
    \item Wird bei jedem Aufruf ausgeführt
    \item Immer aktuell
\end{itemize}
\end{column}

\begin{column}{0.5\textwidth}
\textbf{Vorteile:}
\begin{itemize}
    \item Abstraktion komplexer Abfragen
    \item Zugriffskontrolle (nur View freigeben)
    \item Wiederverwendbarkeit
\end{itemize}
\end{column}
\end{columns}
\end{frame}

\begin{frame}[fragile]{Views erstellen und verwalten}
\begin{lstlisting}
-- View erstellen
CREATE VIEW team_statistik AS
SELECT Mannschaft,
       Punkte,
       ToreGeschossen,
       ToreKassiert,
       ROUND(CAST(Siege AS FLOAT) / Spiele * 100, 1) AS Siegquote
FROM bundesliga;

-- View aktualisieren (DROP + CREATE oder:)
CREATE OR REPLACE VIEW team_statistik AS
SELECT ... -- neue Definition

-- View löschen
DROP VIEW team_statistik;

-- View anzeigen (DuckDB)
SHOW TABLES;  -- Views werden auch gelistet
\end{lstlisting}
\end{frame}

\begin{frame}{Views: Wann einsetzen?}
\begin{columns}[T]
\begin{column}{0.5\textwidth}
\textbf{\cmark\ Gute Anwendungsfälle:}
\begin{itemize}
    \item Komplexe JOINs abstrahieren
    \item Berechnete Spalten bereitstellen
    \item Zugriff auf Teilmengen erlauben
    \item Konsistente Schnittstelle für Anwendungen
\end{itemize}
\end{column}

\begin{column}{0.5\textwidth}
\textbf{\xmark\ Weniger geeignet:}
\begin{itemize}
    \item Sehr häufige, teure Abfragen\\
          $\rightarrow$ Materialized View
    \item Einfache Abfragen\\
          $\rightarrow$ Direkt abfragen
    \item Einmalige Analysen\\
          $\rightarrow$ CTEs nutzen
\end{itemize}
\end{column}
\end{columns}

\vspace{1em}

\begin{exampleblock}{Materialized Views (Konzept)}
\textbf{Materialized Views} speichern das Ergebnis physisch -- schneller, aber muss aktualisiert werden. Nicht in allen Datenbanken verfügbar.
\end{exampleblock}
\end{frame}

%%%%%%%%%%%%%%%%%%%%%%%%%%%%%%%%%%%%%%%%%%%%%%%%%%%%%%%%%%%%%%%%%%%%%%%%%%%%%%%%%%%%%
\showagenda{5}
%%%%%%%%%%%%%%%%%%%%%%%%%%%%%%%%%%%%%%%%%%%%%%%%%%%%%%%%%%%%%%%%%%%%%%%%%%%%%%%%%%%%%

\begin{frame}{Transaktionen: Warum?}
\textbf{Problem:} Mehrere zusammengehörige Änderungen

\begin{center}
\begin{tikzpicture}
    \node[problembox, align=center] (p1) at (0,0) {Überweisung:\\100EUR von A nach B};

    \node[sqlbox, minimum width=3cm, align=center] (s1) at (5,1) {UPDATE A\\SET Saldo = Saldo - 100};
    \node[sqlbox, minimum width=3cm, align=center] (s2) at (5,-1) {UPDATE B\\SET Saldo = Saldo + 100};

    \draw[sqlarrow] (p1) -- (s1);
    \draw[sqlarrow] (p1) -- (s2);

    \node[problembox, draw=red, fill=red!20, align=center] (crash) at (10,0) {Absturz\\nach S1?};
    \draw[roadmaparrow, red] (s1) -- (crash);
\end{tikzpicture}
\end{center}

\vspace{0.5em}

\begin{alertblock}{Das Problem}
A hat 100EUR weniger, B hat nicht mehr bekommen!
\end{alertblock}

\textbf{Lösung:} Transaktionen -- ``Alles oder nichts''
\end{frame}

\begin{frame}{ACID-Eigenschaften}
\begin{center}
\begin{tikzpicture}
    \node[acidbox, fill=red!20] (a) at (0,0) {\shortstack{\textbf{A}tomicity\\Alles oder nichts}};
    \node[acidbox, fill=blue!20] (c) at (4,0) {\shortstack{\textbf{C}onsistency\\Konsistenz}};
    \node[acidbox, fill=green!20] (i) at (8,0) {\shortstack{\textbf{I}solation\\Isolation}};
    \node[acidbox, fill=yellow!30] (d) at (12,0) {\shortstack{\textbf{D}urability\\Dauerhaftigkeit}};
\end{tikzpicture}
\end{center}

\vspace{1em}

\begin{description}
    \item[Atomicity:] Transaktion wird ganz oder gar nicht ausgeführt
    \item[Consistency:] Datenbank bleibt in konsistentem Zustand
    \item[Isolation:] Parallele Transaktionen beeinflussen sich nicht
    \item[Durability:] Bestätigte Änderungen überleben Systemausfälle
\end{description}
\end{frame}

\begin{frame}[fragile]{Transaktion: Syntax}
\begin{lstlisting}
BEGIN TRANSACTION;

-- Alle Änderungen hier sind Teil der Transaktion
UPDATE konten SET saldo = saldo - 100 WHERE konto_id = 'A';
UPDATE konten SET saldo = saldo + 100 WHERE konto_id = 'B';

-- Alles erfolgreich? Speichern!
COMMIT;

-- Oder: Fehler? Alles rückgängig!
ROLLBACK;
\end{lstlisting}

\vspace{0.5em}

\begin{columns}[T]
\begin{column}{0.5\textwidth}
\textbf{COMMIT:}
\begin{itemize}
    \item Alle Änderungen werden permanent
    \item Transaktion ist abgeschlossen
\end{itemize}
\end{column}

\begin{column}{0.5\textwidth}
\textbf{ROLLBACK:}
\begin{itemize}
    \item Alle Änderungen werden verworfen
    \item Zustand vor BEGIN wiederhergestellt
\end{itemize}
\end{column}
\end{columns}
\end{frame}

\begin{frame}[fragile]{Beispiel: Sichere Überweisung}
\begin{lstlisting}
BEGIN TRANSACTION;

-- Prüfen: Hat A genug Geld?
SELECT saldo FROM konten WHERE konto_id = 'A';
-- Annahme: 500EUR -> OK

-- Abbuchen
UPDATE konten SET saldo = saldo - 100 WHERE konto_id = 'A';

-- Gutschreiben
UPDATE konten SET saldo = saldo + 100 WHERE konto_id = 'B';

-- Prüfen: Saldo von A nicht negativ?
SELECT saldo FROM konten WHERE konto_id = 'A';
-- Falls < 0: ROLLBACK; sonst:

COMMIT;
\end{lstlisting}

\textbf{Bei Fehler oder Absturz:} Automatisches ROLLBACK (Atomicity)
\end{frame}

%%%%%%%%%%%%%%%%%%%%%%%%%%%%%%%%%%%%%%%%%%%%%%%%%%%%%%%%%%%%%%%%%%%%%%%%%%%%%%%%%%%%%
\showagenda{6}
%%%%%%%%%%%%%%%%%%%%%%%%%%%%%%%%%%%%%%%%%%%%%%%%%%%%%%%%%%%%%%%%%%%%%%%%%%%%%%%%%%%%%

\begin{frame}{Concurrency-Probleme}
\textbf{Was passiert bei gleichzeitigen Zugriffen?}

\begin{center}
\begin{tikzpicture}[scale=0.9]
    % Timeline
    \draw[->, thick] (0,0) -- (10,0) node[right] {Zeit};

    % Transaction 1
    \node[sqlbox, minimum width=1.5cm] (t1a) at (1,1.5) {T1: Read};
    \node[sqlbox, minimum width=1.5cm] (t1b) at (4,1.5) {T1: Write};
    \node[sqlbox, minimum width=1.5cm] (t1c) at (7,1.5) {T1: Commit};
    \draw[thick, IMSBlue] (t1a) -- (t1b) -- (t1c);

    % Transaction 2
    \node[sqlbox, minimum width=1.5cm, fill=IMSOrange!20] (t2a) at (2.5,-1.5) {T2: Read};
    \node[sqlbox, minimum width=1.5cm, fill=IMSOrange!20] (t2b) at (5.5,-1.5) {T2: Write};
    \node[sqlbox, minimum width=1.5cm, fill=IMSOrange!20] (t2c) at (8.5,-1.5) {T2: Commit};
    \draw[thick, IMSOrange] (t2a) -- (t2b) -- (t2c);

    % Labels
    \node at (-1, 1.5) {Transaktion 1};
    \node at (-1, -1.5) {Transaktion 2};
\end{tikzpicture}
\end{center}

\vspace{0.5em}

\textbf{Klassische Probleme:}
\begin{enumerate}
    \item Lost Update
    \item Dirty Read
    \item Non-Repeatable Read
    \item Phantom Read
\end{enumerate}
\end{frame}

\begin{frame}{Lost Update}
\textbf{Zwei Transaktionen überschreiben sich gegenseitig}

\begin{center}
\begin{tikzpicture}[scale=0.85]
    % Initial state
    \node[tablebox, minimum width=2cm] at (0,3) {Saldo: 100EUR};

    % Timeline
    \draw[->, thick] (0,0) -- (11,0) node[right] {Zeit};

    % T1
    \node[sqlbox, font=\tiny] (t1r) at (2,1.5) {Read: 100EUR};
    \node[sqlbox, font=\tiny] (t1w) at (8,1.5) {Write: 150EUR};
    \draw[thick, IMSBlue] (t1r) -- (t1w);
    \node at (-0.5,1.5) {\footnotesize T1 (+50EUR)};

    % T2
    \node[sqlbox, font=\tiny, fill=IMSOrange!20] (t2r) at (4,-1.5) {Read: 100EUR};
    \node[sqlbox, font=\tiny, fill=IMSOrange!20] (t2w) at (6,-1.5) {Write: 120EUR};
    \draw[thick, IMSOrange] (t2r) -- (t2w);
    \node at (-0.5,-1.5) {\footnotesize T2 (+20EUR)};

    % Result
    \node[problembox, font=\small, align=center] at (9.5,-1.5) {Ergebnis: 150 EUR\\statt 170 EUR!};
\end{tikzpicture}
\end{center}

\vspace{0.5em}

\textbf{Problem:} T2's Änderung (+20EUR) geht verloren, weil T1 den alten Wert überschreibt.
\end{frame}

\begin{frame}{Dirty Read}
\textbf{Lesen von nicht-committeten Daten}

\begin{center}
\begin{tikzpicture}[scale=0.85]
    % Timeline
    \draw[->, thick] (0,0) -- (11,0) node[right] {Zeit};

    % T1
    \node[sqlbox, font=\tiny] (t1w) at (2,1.5) {Write: 150EUR};
    \node[sqlbox, font=\tiny, fill=red!20] (t1r) at (7,1.5) {ROLLBACK};
    \draw[thick, IMSBlue] (t1w) -- (t1r);
    \node at (-0.5,1.5) {\footnotesize T1};

    % T2
    \node[sqlbox, font=\tiny, fill=IMSOrange!20] (t2r) at (4,-1.5) {Read: 150EUR};
    \node[sqlbox, font=\tiny, fill=IMSOrange!20] (t2c) at (9,-1.5) {Commit};
    \draw[thick, IMSOrange] (t2r) -- (t2c);
    \node at (-0.5,-1.5) {\footnotesize T2};

    % Problem
    \draw[dashed, red, thick] (t1w) -- (t2r);
    \node[problembox, font=\small, align=center] at (4,-3) {T2 sieht Wert,\\der nie existierte!};
\end{tikzpicture}
\end{center}

\vspace{0.5em}

\textbf{Problem:} T2 arbeitet mit Daten, die T1 später zurückrollt.
\end{frame}

\begin{frame}{Isolation Levels}
\textbf{Trade-off:} Mehr Isolation = weniger Parallelität

\begin{center}
\footnotesize
\begin{tabular}{|l|c|c|c|c|}
\hline
\rowcolor{IMSBlue!20}
\textbf{Isolation Level} & \textbf{Dirty Read} & \textbf{Non-Rep. Read} & \textbf{Phantom} & \textbf{Performance} \\
\hline
READ UNCOMMITTED & \xmark & \xmark & \xmark & ⚡⚡⚡ \\
\hline
READ COMMITTED & \cmark & \xmark & \xmark & ⚡⚡ \\
\hline
REPEATABLE READ & \cmark & \cmark & \xmark & ⚡ \\
\hline
SERIALIZABLE & \cmark & \cmark & \cmark & 🐢 \\
\hline
\end{tabular}
\end{center}

\vspace{0.5em}

\begin{itemize}
    \item \textbf{READ COMMITTED} ist oft der Standard
    \item \textbf{SERIALIZABLE} verhält sich wie sequentielle Ausführung
    \item DuckDB: Verwendet optimistisches Locking
\end{itemize}
\end{frame}

%%%%%%%%%%%%%%%%%%%%%%%%%%%%%%%%%%%%%%%%%%%%%%%%%%%%%%%%%%%%%%%%%%%%%%%%%%%%%%%%%%%%%
% HANDS-ON 3
{
\setbeamercolor{background canvas}{bg=IMSOrange!15}
\begin{frame}[plain]
\vfill
\begin{center}
{\Huge\color{IMSOrange} Hands-on}\\[1em]
{\Large Views \& Transaktionen}\\[2em]
{\large\ttfamily marimo: 10-subqueries-views-transaktionen.py}\\[1em]
{\normalsize Aufgaben 10.6 -- 10.8}
\end{center}
\vfill
\end{frame}
}

%%%%%%%%%%%%%%%%%%%%%%%%%%%%%%%%%%%%%%%%%%%%%%%%%%%%%%%%%%%%%%%%%%%%%%%%%%%%%%%%%%%%%
\showagenda{7}
%%%%%%%%%%%%%%%%%%%%%%%%%%%%%%%%%%%%%%%%%%%%%%%%%%%%%%%%%%%%%%%%%%%%%%%%%%%%%%%%%%%%%

\begin{frame}{Zusammenfassung}
\begin{columns}[T]
\begin{column}{0.5\textwidth}
\textbf{Subqueries:}
\begin{itemize}
    \item Scalar: Ein Wert
    \item Column: Liste (IN, ANY, ALL)
    \item Table: Zwischentabelle
    \item Korreliert: Referenz auf äußere Abfrage
    \item EXISTS: Existenzprüfung
\end{itemize}

\vspace{0.5em}

\textbf{CTEs (WITH):}
\begin{itemize}
    \item Benannte Zwischenergebnisse
    \item Mehrfach verwendbar
    \item Rekursiv möglich
\end{itemize}
\end{column}

\begin{column}{0.5\textwidth}
\textbf{Views:}
\begin{itemize}
    \item Gespeicherte Abfragen
    \item Verhält sich wie Tabelle
    \item Abstraktion \& Zugriffskontrolle
\end{itemize}

\vspace{0.5em}

\textbf{Transaktionen:}
\begin{itemize}
    \item ACID-Eigenschaften
    \item BEGIN, COMMIT, ROLLBACK
    \item Concurrency-Probleme
    \item Isolation Levels
\end{itemize}
\end{column}
\end{columns}
\end{frame}

\begin{frame}{Entscheidungshilfe: Was wann nutzen?}
\begin{center}
\begin{tikzpicture}[
    decision/.style={diamond, draw=IMSBlue, fill=IMSBlue!10, thick,
        aspect=2, minimum width=2.5cm, font=\small},
    result/.style={rectangle, rounded corners, draw=IMSOrange, fill=IMSOrange!10,
        thick, minimum width=2cm, minimum height=0.8cm, font=\small\bfseries}
]
    \node[decision] (q1) at (0,0) {Einmalig?};
    \node[decision] (q2) at (-3,-2) {Komplex?};
    \node[decision, align=center] (q3) at (3,-2) {Mehrfach\\in Session?};

    \node[result] (r1) at (-5,-4) {Subquery};
    \node[result] (r2) at (-1,-4) {CTE};
    \node[result] (r3) at (3,-4) {View};

    \draw[->, thick] (q1) -- node[left, font=\footnotesize] {ja} (q2);
    \draw[->, thick] (q1) -- node[right, font=\footnotesize] {nein} (q3);
    \draw[->, thick] (q2) -- node[left, font=\footnotesize] {nein} (r1);
    \draw[->, thick] (q2) -- node[right, font=\footnotesize] {ja} (r2);
    \draw[->, thick] (q3) -- node[right, font=\footnotesize] {ja/nein} (r3);
\end{tikzpicture}
\end{center}
\end{frame}

\begin{frame}{Ausblick: Vorlesung 11}
\textbf{Explorative Datenanalyse (EDA)}

\begin{itemize}
    \item Systematische Datenexploration
    \item Univariate \& bivariate Analyse
    \item Visualisierung mit SQL + Python
    \item CASE WHEN für Kategorisierung
    \item Binning und Histogramme
\end{itemize}

\vspace{1em}

\begin{exampleblock}{Vorbereitung}
\begin{itemize}
    \item Notebook 10 abschließen
    \item Wiederholen: Aggregatfunktionen, GROUP BY
\end{itemize}
\end{exampleblock}
\end{frame}

\begin{frame}[plain]
\vfill
\begin{center}
{\Huge Fragen?}\\[2em]
{\Large\ttfamily christoph.flath@uni-wuerzburg.de}
\end{center}
\vfill
\end{frame}

\end{document}
